\section{Related Work}
\label{sec:related}

% \noindent \textbf{Long-Tailed Learning.}
% % Long-tailed learning mainly seeks to address the issue of performance bias toward head classes introduced by class imbalance.
% Long-tailed learning primarily aims to mitigate the performance bias toward head classes arising from class imbalance. Earlier approaches that train models from scratch typically focus on data re-sampling, loss re-weighting, or classifier calibration. These methods aims to emphasize tail classes to alleviate the effects of bias. For instance, logit adjustment (LA) leverages class priors to raise tail-class logits, achieving a theoretically balanced loss. Weight balancing regularizes excessively large head-class norms to re-balance the classifier. Recently, foundation models have shown remarkable effectiveness in LTL by leveraging their rich pre-trained knowledge to enhance long-tailed performance. A common approach is to combine traditional LTL strategies with fine-tuning. For instance, LPT integrates multiple long-tailed tricks (e.g., re-weighting and decoupling) with prompt tuning. Similarly, LIFT adopts the balanced LA loss under the lightweight fine-tuning paradigm. The central goal of these methods is to adapt upstream knowledge to downstream tasks while mitigating the bias introduced by class imbalance. However, beyond data-introduced bias, we observe that fine-tuning itself introduces a new form of bias, i.e., tail-class samples are easily misclassified after fine-tuning, which has been largely overlooked. This paper, in contrast, aims to leverage the strong priors within foundation models to alleviate such bias.

\noindent \textbf{Long-Tailed Learning with Foundation Models.}
Foundation models have recently demonstrated remarkable effectiveness in LTL by leveraging their rich pre-trained knowledge to enhance long-tailed performance~\cite{tian2022vl, dong2023lpt, wang2024exploring, shi2025lift+,zhang2024long,luo2025long,lt2026hong}. A common approach is to combine various LTL strategies~\cite{cui2019class,zhang2021distribution, menon2021longtail,hong2021disentangling,shi2023re} with fine-tuning. For instance, LPT~\cite{dong2023lpt} integrates multiple long-tailed tricks (e.g., re-weighting, decoupling, and cosine classifier) with prompt tuning~\cite{jia2022visual}. Similarly, LIFT~\cite{shi2024long} adopts the balanced LA loss~\cite{menon2021longtail} under the lightweight fine-tuning paradigm. The central goal of these methods is to adapt upstream knowledge to downstream tasks while mitigating the bias introduced by class imbalance. However, beyond data-introduced bias, we observe that fine-tuning itself introduces \textit{concept confusion}, i.e., tail-class samples are easily confused with semantically related categories after fine-tuning, which has been largely overlooked. This paper, in contrast, aims to leverage the strong priors within foundation models to alleviate such bias.
% However, distinct from bias incurred by data, this paper reveals that fine-tuning introduces a new form of bias (even with the advanced LA loss): many samples that were correctly classified by zero-shot foundation models become misclassified after fine-tuning, particularly for tail classes.
% It can be specifically divided into two categories: (1) Directly fine-tuning on foundation models. (2) Introducing additional priors from foundation models.
% \vspace{0.5em}

\noindent \textbf{Pseudo-Labeling.}
% Pseudo-labeling mainly aims to generate artificial labels for unlabeled data and train the models in a self-training manner. For instance, FixMatch employs a fixed confidence threshold to select high-quality pseudo-labels. Similarly, SoftMatch adopts a soft threshold to exploit both high-quality and high-quantity pseudo-labels. These approaches, specifically, aim to generate a single pseudo-label for each unlabeled data to harness the unlabeled samples. Distinct from these approaches, this paper aims to leverage the strong priors of foundation models to generate multiple pseudo-labels, aiming to fully exploit the labeled samples to better learn inter-class correlation.
Pseudo-labeling primarily aims to generate artificial labels for unlabeled samples to better exploit unlabeled data~\cite{hu2021simple, wang2022debiased, abdelfattah2023cdul}. For instance, FixMatch~\cite{sohn2020fixmatch} uses a fixed confidence threshold to select high-quality pseudo-labels for self-training. SoftMatch~\cite{chen2023softmatch}, on the other hand, employs a soft threshold to balance quality and quantity of pseudo-labels. OTAMatch~\cite{zhang2024otamatch} further advances this paradigm by formulating pseudo-labeling as an optimal transport assignment process. The central goal of these approaches is to generate one pseudo-label for each unlabeled sample for supervised learning. In contrast, existing pseudo-labeling methods do not focus on long-tail learning, and they aim to infer a single ground-truth label, whereas we instead seek multiple labels to identify semantically related categories.
% \vspace{0.5em}

\noindent \textbf{Parameter-Efficient Fine-Tuning.}
This paper mainly focuses on parameter-efficient fine-tuning (PEFT)~\cite{xin2024parameter, mai2025lessons, fu2025dtl}, as Shi \textit{et al.}~\cite{shi2024long} demonstrate that full fine-tuning can significantly degrade performance on tail classes, whereas lightweight fine-tuning better preserves upstream knowledge and benefits them. Specifically, PEFT approaches can be grouped into three categories: (i) reparameter tuning (e.g., LoRA~\cite{hu2022lora}), which integrates newly trainable parameters into the pre-trained weights; (ii) prompt tuning (e.g., VPT~\cite{jia2022visual}), which injects learnable prompt tokens into the input space of Transformer blocks; and (iii) adapter tuning (e.g., Adapter~\cite{houlsby2019parameter} and Adaptformer~\cite{chen2022adaptformer}), which inserts extra learnable MLP modules within each Transformer block.
